\documentclass[11pt]{article}

\usepackage[margin=1in]{geometry}
\usepackage{amsmath}
\usepackage{booktabs}
\usepackage{longtable}
\usepackage{array}
\usepackage[T1]{fontenc}
\usepackage{lmodern}

\title{Downstream Metrics Report}
\author{Direction 1 Common Benchmark}
\date{}

\begin{document}
\maketitle

\section{Purpose}

This report explains the downstream metrics in
\texttt{outputs/direction\_1\_run/downstream\_metrics\_aggregated.csv}.
Downstream metrics evaluate whether a reduced representation is still useful
for a later supervised classification task.

For each dataset and method, the benchmark fits a dimensionality-reduction
method on the training data, transforms the train and test data, trains a
classifier on the reduced training data, and evaluates predictions on the
reduced test data. The \texttt{\_mean} columns are averages across seeds, and
the \texttt{\_std} columns measure variation across seeds.

\section{Metric Definitions}

\begin{longtable}{p{0.22\linewidth} p{0.58\linewidth} p{0.14\linewidth}}
\toprule
\textbf{Metric} & \textbf{How it is calculated} & \textbf{Better value} \\
\midrule
\endhead

\texttt{accuracy}
&
Accuracy is the fraction of test examples whose predicted class exactly
matches the true class:
\[
\mathrm{accuracy} =
\frac{\#\{\hat{y}_i = y_i\}}{N}.
\]
Here $N$ is the number of test examples, $y_i$ is the true class, and
$\hat{y}_i$ is the predicted class.
&
Higher
\\
\addlinespace

\texttt{macro\_f1}
&
For each class, compute precision, recall, and F1:
\[
\mathrm{F1}_c =
\frac{2 \cdot \mathrm{precision}_c \cdot \mathrm{recall}_c}
{\mathrm{precision}_c + \mathrm{recall}_c}.
\]
Then average F1 equally across classes:
\[
\mathrm{macro\_F1} = \frac{1}{C}\sum_{c=1}^{C}\mathrm{F1}_c.
\]
Macro F1 is useful when class balance matters because every class receives
equal weight.
&
Higher
\\
\addlinespace

\texttt{elapsed\_seconds}
&
Runtime for fitting the reduction method, transforming the data, and evaluating
the downstream classifier:
\[
t_{\mathrm{elapsed}} = t_{\mathrm{finish}} - t_{\mathrm{start}}.
\]
This measures computational cost, not predictive quality.
&
Lower
\\

\bottomrule
\end{longtable}

\section{Best Method Per Dataset}

The following table selects the best method using the aggregated mean values.
For accuracy and macro F1, the highest mean is best. For elapsed time, the
lowest mean is best.

\begin{longtable}{p{0.22\linewidth} p{0.20\linewidth} p{0.12\linewidth} p{0.30\linewidth} r}
\toprule
\textbf{Dataset} & \textbf{Metric} & \textbf{Direction} & \textbf{Best method} & \textbf{Mean} \\
\midrule
\endhead

\texttt{digits\_clean} & \texttt{accuracy} & Higher & \texttt{pca\_32} & 0.9578 \\
\texttt{digits\_clean} & \texttt{macro\_f1} & Higher & \texttt{pca\_32} & 0.9581 \\
\texttt{digits\_clean} & \texttt{elapsed\_seconds} & Lower & \texttt{pca\_32} & 0.0302 \\
\addlinespace

\texttt{digits\_noisy} & \texttt{accuracy} & Higher & \texttt{pca\_32} & 0.9478 \\
\texttt{digits\_noisy} & \texttt{macro\_f1} & Higher & \texttt{pca\_32} & 0.9480 \\
\texttt{digits\_noisy} & \texttt{elapsed\_seconds} & Lower & \texttt{pca\_32} & 0.0326 \\
\addlinespace

\texttt{mnist\_local} & \texttt{accuracy} & Higher & \texttt{cnn\_autoencoder\_32} & 0.8817 \\
\texttt{mnist\_local} & \texttt{macro\_f1} & Higher & \texttt{cnn\_autoencoder\_32} & 0.8812 \\
\texttt{mnist\_local} & \texttt{elapsed\_seconds} & Lower & \texttt{pca\_32} & 0.0436 \\
\addlinespace

\texttt{twenty\_newsgroups} & \texttt{accuracy} & Higher & \texttt{pca\_32} & 0.8542 \\
\texttt{twenty\_newsgroups} & \texttt{macro\_f1} & Higher & \texttt{pca\_32} & 0.8534 \\
\texttt{twenty\_newsgroups} & \texttt{elapsed\_seconds} & Lower & \texttt{gaussian\_random\_projection\_32} & 0.0078 \\

\bottomrule
\end{longtable}

\section{Interpretation}

The downstream task asks whether reduced features remain useful for
classification. In these results, \texttt{pca\_32} performs best on
\texttt{digits\_clean}, \texttt{digits\_noisy}, and
\texttt{twenty\_newsgroups} for predictive accuracy and macro F1. On
\texttt{mnist\_local}, \texttt{cnn\_autoencoder\_32} gives the best predictive
scores, while \texttt{pca\_32} is fastest.

The runtime winners are not always the predictive winners. Therefore, if the
goal is accuracy, compare \texttt{accuracy} and \texttt{macro\_f1}; if the goal
is speed, compare \texttt{elapsed\_seconds}.

\end{document}
